\documentclass[10pt,a4paper]{amsart}
\usepackage[margin=19mm]{geometry}
\usepackage{amsmath,amssymb,mathtools}
\usepackage[scr=rsfs,cal=euler]{mathalfa}
\usepackage{xfrac}
\usepackage[unicode,hidelinks,bookmarks=false]{hyperref}
\newcommand{\ie}{i.e.\ }
\newcommand{\cf}{cf.\ }
\newcommand{\resp}{respectively\ }
\newcommand{\Subsection}{\subsection}
\providecommand{\nobreakdash}{\nobreak\hbox{-}}
\setlength{\emergencystretch}{2em}
\multlinegap0pt
\usepackage{bm}
\usepackage{bbm}
\usepackage{dsfont}
\usepackage{xspace}
\usepackage{cancel}
\usepackage{mathtools}
\usepackage{amsthm}
\makeatletter
\@ifundefined{proposition}{\newtheorem{proposition}{Proposition}}{}
\makeatother
\title[Making Graphs Irregular through Irregularising Walks]{Making Graphs Irregular through Irregularising Walks}
\author{Julien Bensmail, Romain Bourneuf, Paul Colinot, Samuel Humeau, Timoth\'ee Martinod}
\date{Source: arXiv:2506.21254v1 (CC BY 4.0)}
\begin{document}
\maketitle

\noindent\emph{Source and licence.} Making Graphs Irregular through Irregularising Walks. Version arXiv:2506.21254v1 (CC BY 4.0), distributed under the Creative Commons Attribution 4.0 International licence, as recorded in the arXiv licence metadata for this exact version and shown on the abstract page https://arxiv.org/abs/2506.21254v1. Full rights notice: licenses/arxiv-2506.21254-CC-BY-4.0.txt.\par\smallskip

\noindent\textit{Editorial scope.} Counted unit: the Proposition whose source label is \texttt{proposition:walk\_shapes} in the article (source0.tex lines 878--881, source label \texttt{proposition:walk\_shapes}), delivered together with the article's own complete original proof of it (source0.tex lines 882--917, one \texttt{proof} environment of 36 lines). No mathematics is written, repaired or re-derived: statement and proof are the article's own text line for line. Required context retained uncounted: none. Every definition, display and label of those lines is retained unchanged. Omitted from this package and not redistributed: 1--877, 918--2069. Nothing inside a retained excerpt is omitted or altered: every retained body is delivered exactly as the article's lines are written, so no omission receipt of any kind accompanies this package. Citations: the retained excerpts carry no citation command at all, so no bibliography entry of the article is transcribed and no reference is redistributed. Reference adaptation: none was needed and none was made; every reference inside the delivered unit resolves inside it, so no unresolved reference or citation is produced. Printed slips kept literal: no printed word, symbol, hypothesis or number of the counted statement or of its proof was altered. Layout and class: the article's own class is \texttt{elsarticle} with packages including fontenc, inputenc, babel, amsmath, amsthm, enumerate, amssymb, amscd, vmargin, url, xcolor, stmaryrd, tikz, pgf; the wrapper is the lane's pinned \texttt{amsart} editorial wrapper at 10pt on A4, which restates the theorem-like environments the retained text needs and adds the article's own macro definitions that the retained text needs (none), with extra wrapper packages (bm, bbm, dsfont, xspace, cancel, mathtools). The delivered numbering is the unit's own. Assets: no retained span contains a figure environment, a graphics-inclusion command, a PDF-page inclusion, a table or a TikZ picture; no image file is copied into this package, the package declares no asset dependency and no image file is redistributed with it. Licence: version arXiv:2506.21254v1 of this article is distributed under the Creative Commons Attribution 4.0 International licence, as recorded in the arXiv licence metadata for this exact version and shown on the abstract page https://arxiv.org/abs/2506.21254v1. A full rights notice is delivered at licenses/arxiv-2506.21254-CC-BY-4.0.txt. Characters: every retained excerpt is pure ASCII, so the delivered engine, XeLaTeX with its default Latin Modern text font, reports no missing character.
\begin{proposition}\label{proposition:walk_shapes}
  Let $G$ be a graph, and $W$ be a walk of $G$. Denote by $E_o$ and $E_e$ the sets of edges traversed by $W$
  an odd number of times, and a non-zero even number of times, respectively. Then, $W$ is equivalent to a walk $W_S$ that follows a walk $S$ in which each edge of $E_o$ appears once, and each edge of $E_e$ appears at most twice.
\end{proposition}

\begin{proof}
  We replace $W$ by an equivalent walk to gather occurrences of each edge.
  
  First, we show how this is done for individual edges.
  Let $e=xy$ be an edge traversed by $W$. Assume that $e$ occurs $n\ge 1$
  times in $W$ and let $W_1,\dots,W_{n+1}$ be the largest subwalks of $W$ not containing $e$; that is,
  \[
    W=W_1e_1W_2e_2\dots W_ne_nW_{n+1},
  \]
  with $e_i$ being oriented edges in $\{xy,yx\}$. Up to reversing them, the subwalks $W_i$ for $i\notin\{1,n+1\}$ can be partitioned into three sets:
  the ones in $\mathcal{W}_x$ going from $x$ to $x$, the ones in $\mathcal{W}_y$ going from $y$ to $y$, and those in $\mathcal{W}_{xy}$
  going from $x$ to $y$. W.l.o.g., suppose that $W_1$ goes from the first vertex of $W$ to $x$. The walk $W_{n+1}$ can start either from $x$ or from $y$, and then goes on to the last vertex of $W$. These two cases depend on the parity of $n+|\mathcal W_{xy}|$: if this number is even then $W_{n+1}$ must start from $x$, and if it is odd it starts from $y$. This fact will justify that equivalent walks below are well-defined.
  %\textcolor{red}{Rom: You need to handle $W_1$ and $W_{n+1}$ separately I guess. We can consider four cases depending on whether they have the same endpoint in $\{x, y\}$ or not, and whether $n$ is even or not. If yes and yes, say the endpoint is $x$, we could follow $W_1$ then do $\mathcal{W}_x$ then take $e$ then do $\mathcal{W}_y$ then $\mathcal{W}_{xy}$ then take $e$ again as much as we need, which brings us back to $x$ (for parity reasons) where we can finally take $W_{n+1}$. If yes and no then $W_{xy}$ is odd for parity reasons, so we could follow $W_1$ then do $\mathcal{W}_x$ then take $e$ as much as we need, which brings us to $y$, then do $\mathcal{W}_y$ then $\mathcal{W}_{xy}$, which brings us back to $x$ where we can finally take $W_{n+1}$. If no, we could follow $W_1$ then do $\mathcal{W}_x$ then $\mathcal{W}_{xy}$ then take $e$ as much as we need, which brings us to $y$ (for parity reasons) where we can finally do $\mathcal{W}_y$ before taking $W_{n+1}$.}
  
  \begin{itemize}
      \item Assume first that $\mathcal{W}_{xy}$ is not empty. In this case, choose $W'\in\mathcal{W}_{xy}$, and replace
  $W$ by
  the equivalent walk
  \[
      W_1\left( \prod_{W_x\in\mathcal{W}_x}W_x \right)W'\left( \prod_{W_y\in\mathcal{W}_y}W_y \right) e^n\left(\prod_{W_{xy}\in\mathcal{W}_{xy}\backslash\{W'\}}W_{xy}\right)W_{n+1},
  \]
  where products are concatenation of walks, and the $n$ occurences of $e$, but also the walks under the products, are oriented in the right ways for the final walk to be well defined (e.g.~$e^1=yx$, $e^2=yxy$ and not $e^2=yxyx$, etc).
  
  \item Assume now that $\mathcal{W}_{xy}$ is empty. If $n$ is odd, then replace $W$ by the equivalent walk
  \[
    W_1\left( \prod_{W_x\in\mathcal{W}_x}W_x \right)e^n\left( \prod_{W_y\in\mathcal{W}_y}W_y \right)W_{n+1}.
  \]
  If $n$ is even, then replace $W$ by the equivalent walk
  \[
    W_1\left( \prod_{W_x\in\mathcal{W}_x}W_x \right)e\left( \prod_{W_y\in\mathcal{W}_y}W_y \right)e^{n-1}W_{n+1}.
  \]
  \end{itemize}
  
  The fact that the parts $W_i$ of $W$ are not modified by this process allows us to apply it recursively to each edge of $W$.
  At the end, we are left with a walk in which each edge appears in at most two places.
\end{proof}

\end{document}
